\documentclass[10pt,a4paper]{amsart}
\usepackage{amsmath,amssymb,amsthm,graphicx}
% Retained minimal preamble for the Stacks Project chapter varieties.tex.
% Verbatim excerpt of the chapter preamble: Xy-pic setup (native preamble lines 14-18), the native theorem
% declarations with the mathematical macros (native preamble lines 154-217), and the project own no-op
% definitions of the reference/slogan/history display environments (native wrapper.tex lines 1-3) that the
% counted statement uses for its slogan. Copied mechanically; nothing added, removed or redefined.
% For commutative diagrams we use Xy-pic
\usepackage[all]{xy}

% We use 2cell for 2-commutative diagrams.
\xyoption{2cell}
% Theorem environments.
%
\theoremstyle{plain}
\newtheorem{theorem}[subsection]{Theorem}
\newtheorem{proposition}[subsection]{Proposition}
\newtheorem{lemma}[subsection]{Lemma}

\theoremstyle{definition}
\newtheorem{definition}[subsection]{Definition}
\newtheorem{example}[subsection]{Example}
\newtheorem{exercise}[subsection]{Exercise}
\newtheorem{situation}[subsection]{Situation}

\theoremstyle{remark}
\newtheorem{remark}[subsection]{Remark}
\newtheorem{remarks}[subsection]{Remarks}

\numberwithin{equation}{subsection}

% Macros
%
\def\lim{\mathop{\mathrm{lim}}\nolimits}
\def\colim{\mathop{\mathrm{colim}}\nolimits}
\def\Spec{\mathop{\mathrm{Spec}}}
\def\Hom{\mathop{\mathrm{Hom}}\nolimits}
\def\Ext{\mathop{\mathrm{Ext}}\nolimits}
\def\SheafHom{\mathop{\mathcal{H}\!\mathit{om}}\nolimits}
\def\SheafExt{\mathop{\mathcal{E}\!\mathit{xt}}\nolimits}
\def\Sch{\mathit{Sch}}
\def\Mor{\mathop{\mathrm{Mor}}\nolimits}
\def\Ob{\mathop{\mathrm{Ob}}\nolimits}
\def\Sh{\mathop{\mathit{Sh}}\nolimits}
\def\NL{\mathop{N\!L}\nolimits}
\def\CH{\mathop{\mathrm{CH}}\nolimits}
\def\proetale{{pro\text{-}\acute{e}tale}}
\def\etale{{\acute{e}tale}}
\def\QCoh{\mathit{QCoh}}
\def\Ker{\mathop{\mathrm{Ker}}}
\def\Im{\mathop{\mathrm{Im}}}
\def\Coker{\mathop{\mathrm{Coker}}}
\def\Coim{\mathop{\mathrm{Coim}}}

% Boxtimes
%
\DeclareMathSymbol{\boxtimes}{\mathbin}{AMSa}{"02}

%
% Macros for moduli stacks/spaces
%
\def\QCohstack{\mathcal{QC}\!\mathit{oh}}
\def\Cohstack{\mathcal{C}\!\mathit{oh}}
\def\Spacesstack{\mathcal{S}\!\mathit{paces}}
\def\Quotfunctor{\mathrm{Quot}}
\def\Hilbfunctor{\mathrm{Hilb}}
\def\Curvesstack{\mathcal{C}\!\mathit{urves}}
\def\Polarizedstack{\mathcal{P}\!\mathit{olarized}}
\def\Complexesstack{\mathcal{C}\!\mathit{omplexes}}
% \Pic is the operator that assigns to X its picard group, usage \Pic(X)
% \Picardstack_{X/B} denotes the Picard stack of X over B
% \Picardfunctor_{X/B} denotes the Picard functor of X over B
\def\Pic{\mathop{\mathrm{Pic}}\nolimits}
\def\Picardstack{\mathcal{P}\!\mathit{ic}}
\def\Picardfunctor{\mathrm{Pic}}
\def\Deformationcategory{\mathcal{D}\!\mathit{ef}}
\newenvironment{reference}{}{}
\newenvironment{slogan}{}{}
\newenvironment{history}{}{}
\usepackage{cleveref}
\crefname{originallemma}{original lemma}{original lemmas}
\Crefname{originallemma}{Original lemma}{Original lemmas}
\makeatletter
\newlabel{lemma-general-degree-g-line-bundle}{{0B8Z}{source}{}{sourcepointer-1}{}}
\newlabel{lemma-general-degree-g-line-bundle@cref}{{[originallemma][0][]0B8Z}{[1][0][]source}{}{}{}}
\makeatother
\title{Vanishing of the first cohomology of invertible sheaves of degree at least twice the genus minus one on a proper curve}
\author{The Stacks Project Authors}
\date{stacks-project:varieties.tex:local-snapshot-20261011; GNU FDL 1.2 or later}
\begin{document}
\maketitle
\noindent\textit{Editorial note.} All delivered mathematical text is the author's own, copied verbatim from the pinned Stacks Project chapter file varieties.tex; no statement or proof step was rewritten, reordered or removed.\par\medskip
\noindent\textit{Editorial note.} The counted span cites one cross-chapter label, recorded in assembly.external\_references with its canonical tag as a navigation pointer.\par\medskip
\noindent\textit{Editorial note.} The proof cites no bibliography item inside the counted span.\par\medskip
\noindent\textit{Allowed standard prerequisites (editorial).}\begin{itemize}
\item Existence of an invertible sheaf of prescribed degree with a vanishing first cohomology. Stacks Project, Varieties, Lemma lemma-general-degree-g-line-bundle (tag 0B4A). Supplies the auxiliary invertible module N of degree g - 1.
\end{itemize}\medskip
\ifdefined\hypertarget\hypertarget{sourcepointer-1}{}\fi
\noindent\textit{Source reference.} Original lemma 0B8Z: sources/stacks-chow/tags.txt, source lines 14157--14157. Cross-chapter standard prerequisite cited by the author inside the counted span; it is used as a pointer and its own proof is not part of this item.\par\medskip
\begin{lemma}
\label{lemma-vanishing-degree-2g-and-1-line-bundle}
Let $k$ be an algebraically closed field. Let $X$ be a proper curve over $k$.
Set $g = \dim_k H^1(X, \mathcal{O}_X)$. For every invertible
$\mathcal{O}_X$-module $\mathcal{L}$ with $\deg(\mathcal{L}) \geq 2g - 1$
we have $H^1(X, \mathcal{L}) = 0$.
\end{lemma}

\begin{proof}
Let $\mathcal{N}$ be the invertible module we found in
Lemma \ref{lemma-general-degree-g-line-bundle} part (2).
The degree of $\mathcal{N}$ is
$\chi(X, \mathcal{N}) - \chi(X, \mathcal{O}_X) = 0 - (1 - g) = g - 1$.
Hence the degree of $\mathcal{L} \otimes \mathcal{N}^{\otimes - 1}$
is $\deg(\mathcal{L}) - (g - 1) \geq g$.
Hence
$\chi(X, \mathcal{L} \otimes \mathcal{N}^{\otimes -1}) \geq g + 1 - g = 1$.
Thus there is a nonzero global section $s$ whose zero scheme is an
effective Cartier divisor $D$ of degree $\deg(\mathcal{L}) - (g - 1)$.
This gives a short exact sequence
$$
0 \to \mathcal{N} \xrightarrow{s} \mathcal{L} \to i_*(\mathcal{L}|_D) \to 0
$$
where $i : D \to X$ is the inclusion morphism. We conclude that
$H^0(X, \mathcal{L})$ maps isomorphically to $H^0(D, \mathcal{L}|_D)$
which has dimension $\deg(\mathcal{L}) - (g - 1)$. The result follows
from the definition of degree.
\end{proof}

\end{document}
